% Generated draft: requires mathematical and source review.
\documentclass[11pt]{article}
\usepackage[margin=25mm]{geometry}
\usepackage{fontspec}
\setmainfont{Latin Modern Roman}
\usepackage{amsmath,amssymb,mathtools,unicode-math}
\setmathfont{Latin Modern Math}
\usepackage{enumitem,needspace,xurl,hyperref}
\setlength{\parindent}{0pt}
\setlength{\parskip}{5pt}
\setlength{\emergencystretch}{3em}
\newcommand{\N}{\mathbb{N}}
\newcommand{\Z}{\mathbb{Z}}
\newcommand{\Q}{\mathbb{Q}}
\newcommand{\R}{\mathbb{R}}
\newcommand{\C}{\mathbb{C}}
\newcommand{\F}{\mathbb{F}}
\newcommand{\A}{\mathbb{A}}
\newcommand{\PP}{\mathbb{P}}
\newcommand{\E}{\mathbb{E}}
\newcommand{\eps}{\varepsilon}
\newcommand{\dd}{\,\mathrm{d}}
\begin{document}
\section*{UnsolvedMath ID: 701654}
\textbf{Source problem code:} catalogue entry
\subsection{Definitions and mathematical statement}
A group \(G\) is uniformly locally finite if there exists a function \(f\colon \Z_{\geq 1}\to\Z_{\geq 1}\) such that, for every \(n\geq 1\) and every \(g_1,\ldots,g_n\in G\),
\[
\bigl|\langle g_1,\ldots,g_n\rangle\bigr|\leq f(n).
\]
A subgroup is definable relative to a first-order language and a convention about permitted parameters if it is the set of elements satisfying a formula in that language with those permitted parameters. Let \(\mathcal D(G)\) denote the family of subgroups intended by the source's use of ``definable.'' The supplied record does not specify the language, an ambient structure, or whether parameters are allowed; this qualification remains unresolved.

The group \(G\) is nilpotent-by-finite if it has a normal subgroup \(L\trianglelefteq G\) such that \([G:L]<\infty\) and \(L\) is nilpotent. Here nilpotence means that the lower central series, defined by \(\gamma_1(L)=L\) and \(\gamma_{i+1}(L)=[\gamma_i(L),L]\), reaches the trivial subgroup after finitely many steps. For subgroups \(A,B\), the subgroup \([A,B]\) is generated by all \(a^{-1}b^{-1}ab\) with \(a\in A\) and \(b\in B\). The group \(G\) is finite-by-abelian-by-finite if there exist \(G_0\trianglelefteq G\) and \(N\trianglelefteq G_0\) with \([G:G_0]<\infty\), \(|N|<\infty\), and \(G_0/N\) abelian.

For the intended definability convention, suppose that \(G\) is uniformly locally finite and
\[
\forall H,K\in\mathcal D(G),\qquad
[H:H\cap K]<\infty\quad\text{or}\quad[K:H\cap K]<\infty.
\]
Must \(G\) be nilpotent-by-finite? Does it even follow that \(G\) is finite-by-abelian-by-finite? These are the two conclusions requested by the source; a fully specified formulation requires clarification of definability.
\subsection{Short English statement}
F. Wagner asks whether a uniformly locally finite group in which the intersection of any two definable subgroups has finite index in at least one of them must be nilpotent-by-finite, or even finite-by-abelian-by-finite. The meaning of definability requires clarification. The supplied registry labels the problem open; this draft does not independently verify that status.
\subsection{Sources}
\begin{itemize}
\item Local UnsolvedMath registry: supplied complete source record, permanent ID \(701654\), exact source code \texttt{catalogue entry}, title \emph{Kourovka Notebook Problem 19.106}. The record attributes the problem to F. Wagner, with proposed year 2018, and labels it \texttt{open}. Its creation and update timestamps are both \texttt{2026-10-16T14:00:00Z}; these are registry metadata, not independent literature verification.
\item Supplied source attribution: \emph{The Kourovka Notebook (Unsolved Problems in Group Theory)}, 19th issue (2018), original problem 19.106; also identified in the record with the 21st issue edition (2026), \texttt{arXiv:1401.0300v48}. The supplied collection description names E. I. Khukhro and V. D. Mazurov as editors. This bibliographic attribution is supplied evidence and has not been independently checked.
\end{itemize}
\end{document}
